\chapter{存储器件结构：从单元到阵列到接口}

CPU 寄存器堆里一个 bit 由两个交叉耦合反相器保存，主存里一个 bit 由一颗电容保存，闪存里一个 bit 由被困在绝缘层中的电子保存，硬盘里一个 bit 由一小块磁性薄膜的磁化方向保存。同样是 0 和 1，物理载体完全不同，读出方式、写入代价、保持时间、访问粒度和失效模式也完全不同。本章把这四类存储器件的内部结构讲到可以画截面图、算噪声裕量、读时序参数的程度。

\begin{keyidea}
存储层次不是"快慢不同的同一种东西"，而是\hwkey{物理机制完全不同的器件}。SRAM 靠反馈锁存，DRAM 靠电容电荷，NAND 靠绝缘层中的被困电子，HDD 靠磁畴方向。每一层的访问粒度、写入代价和失效模式都由其物理结构决定。
\end{keyidea}

\section{一、SRAM 单元与阵列}

\topic{6T 交叉耦合单元}

SRAM 的基本存储单元由六个 MOSFET 组成：两个交叉耦合的反相器形成双稳态锁存，再加两个访问晶体管把锁存器连接到字线和位线。

\begin{pdfdiagram}[x=1cm,y=1cm]
  % 交叉耦合反相器
  \node[hw state, minimum width=1.2cm, minimum height=0.7cm] (inv1) at (0,0) {P1\\N1};
  \node[hw state, minimum width=1.2cm, minimum height=0.7cm] (inv2) at (3.2,0) {P2\\N2};
  % 存储节点
  \node[BookTeal, font=\small\bfseries] at (-0.9,0) {Q};
  \node[BookTeal, font=\small\bfseries] at (4.1,0) {$\overline{Q}$};
  % 交叉耦合线
  \draw[hw bus] (inv1.east) -- ++(0.3,0) |- (inv2.north west);
  \draw[hw bus] (inv2.west) -- ++(-0.3,0) |- (inv1.south west);
  % 访问晶体管
  \node[hw compute, minimum width=0.9cm, minimum height=0.6cm] (acc1) at (-2.2,0) {M5};
  \node[hw compute, minimum width=0.9cm, minimum height=0.6cm] (acc2) at (5.4,0) {M6};
  % 位线
  \draw[hw bus] (acc1.west) -- ++(-1.0,0) node[left, hw label] {BL};
  \draw[hw bus] (acc2.east) -- ++(1.0,0) node[right, hw label] {$\overline{BL}$};
  % 字线
  \draw[BookAccent, line width=0.7pt] (-2.2,0.7) -- (-2.2,1.3) node[above, hw label] {WL};
  \draw[BookAccent, line width=0.7pt] (5.4,0.7) -- (5.4,1.3) node[above, hw label] {WL};
  \draw[BookAccent, line width=0.5pt, dashed] (-2.2,1.3) -- (5.4,1.3);
  % 电源和地
  \draw[hw return] (inv1.north) -- ++(0,0.5) node[above, hw label] {VDD};
  \draw[hw return] (inv2.north) -- ++(0,0.5) node[above, hw label] {VDD};
  \draw[hw return] (inv1.south) -- ++(0,-0.5) node[below, hw label] {VSS};
  \draw[hw return] (inv2.south) -- ++(0,-0.5) node[below, hw label] {VSS};
  \node[hw label, text width=8cm] at (1.6,-1.6) {6T SRAM 单元：P1/N1 和 P2/N2 交叉耦合形成双稳态；M5/M6 是访问管，栅极接字线 WL，源/漏分别接 Q/$\overline{Q}$ 和位线 BL/$\overline{BL}$。};
\end{pdfdiagram}
\diagramnote{6T SRAM 单元截面逻辑图。WL=0 时 M5/M6 截止，锁存器自保持；WL=1 时 M5/M6 导通，位线可以读出或写入 Q 和 $\overline{Q}$。}

单元的工作原理分三种状态：

\begin{longtable}{L{0.14\linewidth}L{0.36\linewidth}L{0.4\linewidth}}
\toprule
状态 & 字线 WL & 位线行为 \\
\midrule
保持 & 0（M5/M6 截止） & 位线浮空或预充电；锁存器靠反馈自维持 Q 和 $\overline{Q}$ \\
读出 & 1（M5/M6 导通） & 位线预充到 VDD 后，Q 侧通过 M5 微微放电（若 Q=0）或保持高（若 Q=1）；灵敏放大器检测 BL 与 $\overline{BL}$ 的微小压差 \\
写入 & 1（M5/M6 导通） & 外部驱动 BL=0、$\overline{BL}$=1（写 0）或反过来；位线驱动能力必须强过锁存器反馈，才能翻转存储节点 \\
\bottomrule
\end{longtable}

\topic{读操作：电荷分享与灵敏放大}

读出时，位线先被预充到 VDD。字线打开后，若 Q=0（即存储节点为低），M5 导通把 BL 通过 N1 向地放电；若 Q=1，BL 保持高。实际压差极小——典型 6T 单元在 0.9V 工艺下，BL 放电幅度仅 50--100 mV。灵敏放大器（cross-coupled latch 或电流镜）把这个微小差值放大到全摆幅。

\begin{longtable}{L{0.28\linewidth}L{0.22\linewidth}L{0.4\linewidth}}
\toprule
参数 & 典型值（28nm） & 含义 \\
\midrule
BL 预充电压 & VDD (0.9V) & 读前位线被拉到满高 \\
读放电幅度 & 50--100 mV & Q=0 时 BL 下降量 \\
灵敏放大延迟 & 0.3--0.8 ns & 从压差出现到输出全摆幅 \\
读访问时间 & 0.8--2.0 ns & WL 有效到数据输出稳定 \\
\bottomrule
\end{longtable}

\topic{写操作：驱动能力竞争}

写入时，外部驱动 BL 和 $\overline{BL}$ 为互补电平。字线打开后，位线驱动管（通常是列写驱动器中的大尺寸 NMOS）必须把存储节点拉过反相器的翻转阈值，才能让反馈锁存到新的稳定状态。这要求：

\begin{lstlisting}
write driver strength > pull-up PMOS (P1 or P2) strength
access transistor (M5/M6) must pass enough current
\end{lstlisting}

工艺缩小后，PMOS 上拉管和访问管的尺寸比（cell ratio 和 pull-up ratio）成为写裕量的关键设计参数。

\topic{静态噪声裕量（SNM）}

SNM 衡量单元在保持状态下能容忍多大的外部干扰而不翻转。测量方法是画出两个反相器的电压传输特性（VTC），把其中一条镜像叠加，形成蝶形图（butterfly curve）。蝶形图中最大内切正方形的边长就是 SNM。

\begin{pdfdiagram}[x=1cm,y=1cm]
  % 坐标轴
  \draw[BookRule, line width=0.5pt, ->] (0,0) -- (5.5,0) node[right, hw label] {$V_{in}$ (V)};
  \draw[BookRule, line width=0.5pt, ->] (0,0) -- (0,5.5) node[above, hw label] {$V_{out}$ (V)};
  % VTC 曲线 1（正常）
  \draw[BookAccent, line width=0.8pt] (0.3,5.0) .. controls (1.5,5.0) and (2.0,4.5) .. (2.5,2.5) .. controls (3.0,0.5) and (3.5,0.3) .. (5.0,0.3);
  % VTC 曲线 2（镜像）
  \draw[BookTeal, line width=0.8pt] (5.0,5.0) .. controls (3.5,5.0) and (3.0,4.5) .. (2.5,2.5) .. controls (2.0,0.5) and (1.5,0.3) .. (0.3,0.3);
  % SNM 正方形
  \draw[BookAmber, line width=0.7pt] (1.6,1.6) rectangle (3.4,3.4);
  \node[BookAmber, font=\small\bfseries] at (2.5,2.5) {SNM};
  % 标注
  \node[hw label] at (4.8,5.2) {VDD};
  \node[hw label] at (5.2,0.1) {VDD};
  \node[hw label, text width=7cm] at (2.7,-0.9) {蝶形图：两条 VTC 围成的两个"眼"中，最大内切正方形边长 = SNM。SNM 越大，单元抗干扰能力越强。};
\end{pdfdiagram}
\diagramnote{SRAM 蝶形图。工艺缩小后 VDD 降低、阈值电压波动增大，SNM 缩小，这是 SRAM 在先进节点面积不缩反增的主要原因。}

\begin{longtable}{L{0.22\linewidth}L{0.22\linewidth}L{0.46\linewidth}}
\toprule
工艺节点 & VDD & 典型 SNM \\
\midrule
65 nm & 1.0 V & 250--350 mV \\
28 nm & 0.9 V & 150--250 mV \\
7 nm & 0.7 V & 80--150 mV \\
5 nm & 0.65 V & 60--120 mV（需要 assist 电路） \\
\bottomrule
\end{longtable}

在 7nm 及以下，SNM 已经接近工艺波动幅度，必须引入写辅助电路：负位线（negative BL）、字线下压（WL underdrive）、电源下压（VDD collapse）或读辅助（read VDD boost）。

\topic{从单元到阵列}

一个 64 Mbit SRAM 阵列由 $8192 \times 8192$ 个 6T 单元排成矩阵。组织方式：

\begin{longtable}{L{0.22\linewidth}L{0.34\linewidth}L{0.34\linewidth}}
\toprule
结构 & 作用 & 设计约束 \\
\midrule
字线（WL） & 一行所有单元的访问管栅极连在一起 & WL 驱动能力决定行激活速度；行长 $\propto$ 列数 \\
位线对（BL/$\overline{BL}$） & 一列所有单元的存储节点通过访问管连到位线 & 位线电容 $\propto$ 行数；电容越大，读放电越慢 \\
列多路选择器 & 从多条位线中选一条送到灵敏放大器 & 4:1 或 8:1 mux 减少灵敏放大器数量 \\
灵敏放大器 & 检测位线微小压差并放大到全摆幅 & 每列一个或每几列共享 \\
写驱动器 & 驱动位线对为互补电平 & 驱动能力必须超过单元上拉管 \\
行译码器 & 把行地址译成一条 WL 有效 & 地址位宽决定阵列行数 \\
\bottomrule
\end{longtable}

cache 中的 SRAM 阵列还额外嵌入 ECC 位（通常每 64 数据位加 8 校验位）和冗余行列（修复制造缺陷）。L1 cache 要求单周期访问（$<$1 ns），因此阵列小（32--64 KiB）、位线短、灵敏放大器快；L2/L3 容量大但允许多周期访问，阵列更大、位线更长。

\section{二、DRAM 单元、感放与阵列}

\topic{1T1C 单元：用电容存一个 bit}

DRAM 单元只有一个访问晶体管（1T）加一个存储电容（1C）。电容上有电荷代表 1，无电荷（或电荷少）代表 0。

\begin{pdfdiagram}[x=1cm,y=1cm]
  % 位线
  \draw[hw bus] (0,2.0) -- (0,0.5);
  \node[hw label, above] at (0,2.1) {BL};
  % 访问晶体管
  \node[hw compute, minimum width=1.0cm, minimum height=0.7cm] (acc) at (0,0) {M1};
  \draw[hw bus] (0,0.5) -- (acc.north);
  % 字线
  \draw[BookAccent, line width=0.7pt] (acc.west) -- ++(-1.2,0) node[left, hw label] {WL};
  % 存储电容
  \draw[hw bus] (acc.east) -- ++(1.0,0);
  \draw[line width=0.7pt] (1.8,-0.15) -- (1.8,0.15);
  \draw[line width=0.7pt] (2.0,-0.15) -- (2.0,0.15);
  \draw[hw bus] (2.0,0) -- ++(0.8,0) node[right, hw label] {VDD/2};
  \node[BookTeal, font=\small\bfseries] at (1.9,-0.5) {$C_S$};
  % 标注
  \node[hw label, text width=7.5cm] at (1.0,-1.4) {1T1C 单元：WL 打开 M1 后，$C_S$ 与 BL 的寄生电容 $C_{BL}$ 发生电荷共享；共享后 BL 电压偏移极小（几十 mV），必须由灵敏放大器锁存到全摆幅。};
\end{pdfdiagram}
\diagramnote{1T1C DRAM 单元。存储电容一端接访问管，另一端接 VDD/2 参考。读出是破坏性的——电荷共享后必须回写。}

\topic{电荷共享定量}

设存储电容 $C_S = 30$ fF，位线寄生电容 $C_{BL} = 300$ fF（典型 10:1 比值）。预充时 BL 被拉到 VDD/2 = 0.5V。

若 $C_S$ 上存的是 1（电压 = VDD = 1.0V），打开访问管后电荷守恒：

$$V_{BL}' = \frac{C_{BL} \cdot 0.5 + C_S \cdot 1.0}{C_{BL} + C_S} = \frac{300 \times 0.5 + 30 \times 1.0}{330} = 0.545 \text{V}$$

BL 上升 45 mV。若 $C_S$ 存的是 0（电压 = 0V），BL 下降 45 mV 到 0.455V。灵敏放大器必须分辨 $\pm 45$ mV 的差值。

\begin{longtable}{L{0.3\linewidth}L{0.2\linewidth}L{0.4\linewidth}}
\toprule
参数 & 典型值 & 设计影响 \\
\midrule
$C_S$ & 20--30 fF & 越小信号越弱； trench/stack 电容工艺维持 $C_S$ \\
$C_{BL}$ & 200--400 fF & 由列中行数决定；行越多越慢 \\
$C_S / C_{BL}$ & 1:8 -- 1:12 & 比值决定信号幅度 \\
信号幅度 & 40--80 mV & 灵敏放大器输入 offset 必须远小于此 \\
\bottomrule
\end{longtable}

\topic{灵敏放大器：交叉耦合锁存}

DRAM 灵敏放大器本身就是一个交叉耦合 NMOS+PMOS 对（与 SRAM 单元结构相同，但尺寸大得多）。电荷共享后，BL 和 $\overline{BL}$ 出现微小压差；灵敏放大器正反馈把这个压差在 1--2 ns 内放大到全摆幅（0 或 VDD）。放大完成后，数据被锁存在灵敏放大器中——这就是 row buffer（行缓冲）。

\begin{pdfdiagram}[x=1cm,y=1cm]
  % 时间轴
  \draw[BookRule, line width=0.5pt, ->] (0,0) -- (10,0) node[right, hw label] {时间};
  % BL 波形
  \draw[BookTeal, line width=0.8pt] (0.5,2.0) -- (2.0,2.0) -- (2.5,2.15) -- (3.5,2.15) -- (4.0,3.5) -- (9.0,3.5);
  \node[hw label, anchor=east] at (0.4,2.0) {BL};
  \node[hw label] at (2.25,2.4) {+45mV};
  \node[hw label] at (5.5,3.7) {VDD};
  % BLbar 波形
  \draw[BookAccent, line width=0.8pt] (0.5,1.5) -- (2.0,1.5) -- (2.5,1.35) -- (3.5,1.35) -- (4.0,0.3) -- (9.0,0.3);
  \node[hw label, anchor=east] at (0.4,1.5) {$\overline{BL}$};
  \node[hw label] at (2.25,1.1) {−45mV};
  \node[hw label] at (5.5,0.1) {0};
  % 阶段标注
  \draw[BookRule, line width=0.4pt, dashed] (2.0,4.0) -- (2.0,-0.3);
  \draw[BookRule, line width=0.4pt, dashed] (4.0,4.0) -- (4.0,-0.3);
  \node[hw label] at (1.2,4.2) {预充 VDD/2};
  \node[hw label] at (3.0,4.2) {电荷共享};
  \node[hw label] at (6.5,4.2) {感放锁存→全摆幅};
\end{pdfdiagram}
\diagramnote{DRAM 读出过程：预充→字线打开→电荷共享产生 $\pm$45mV→灵敏放大器正反馈放大到全摆幅。整个过程是破坏性的，放大后数据必须回写到单元电容。}

\topic{行列译码与 bank 结构}

一个 8 Gbit DDR4 芯片内部组织为 4 个 bank group，每组 4 个 bank，每 bank 约 512 Mbit。每个 bank 内部是一个巨大的二维阵列：

\begin{longtable}{L{0.24\linewidth}L{0.26\linewidth}L{0.4\linewidth}}
\toprule
层级 & 数量（8Gb DDR4） & 说明 \\
\midrule
Bank group & 4 & 不同 group 间命令间隔更短，提高并行度 \\
Bank / group & 4 & 共 16 bank；同一时刻每 bank 只能有一行打开 \\
Row / bank & 64K--256K & 行地址（RAS）选中一行，整行数据进入 row buffer \\
Column / row & 数千 & 列地址（CAS）从 row buffer 中选 8 或 16 bit 输出 \\
\bottomrule
\end{longtable}

行激活（ACTIVATE）打开一整行到 row buffer，耗时 tRCD（~14 ns）；列读（READ）从 row buffer 取数据，耗时 CL（~14 ns）；行预充（PRECHARGE）关闭当前行，耗时 tRP（~14 ns）。若连续访问同一行，只需付 CL；若换行，必须付 tRP + tRCD + CL。

\topic{刷新：电容漏电的代价}

DRAM 电容通过 PN 结漏电，典型保持时间 32--64 ms。必须在数据丢失前把所有行重新读出并回写。标准做法：每 64 ms 完成全部行的刷新，即每 7.8 $\mu$s 发一次 REFRESH 命令刷新一批行（auto-refresh）。

\begin{longtable}{L{0.28\linewidth}L{0.22\linewidth}L{0.4\linewidth}}
\toprule
刷新参数 & 典型值 & 影响 \\
\midrule
保持时间 & 32--64 ms & 超过此时间数据丢失 \\
刷新间隔 & 7.8 $\mu$s & 64 ms / 8192 行 \\
刷新命令占用 & 约 350 ns & 期间 bank 不可访问 \\
刷新带宽损失 & 3--5\% & 每 7.8 $\mu$s 中有 350 ns 不可用 \\
\bottomrule
\end{longtable}

DDR5 引入 same-bank refresh 和 fine-granularity refresh，允许只刷新部分 bank 而其余 bank 继续服务请求，降低刷新对带宽的影响。

\section{三、DDR 信号、PHY 与训练}

\topic{DDR 信号结构}

DDR（Double Data Rate）SDRAM 在时钟的上升沿和下降沿都传输数据。信号分为三组：

\begin{longtable}{L{0.18\linewidth}L{0.32\linewidth}L{0.4\linewidth}}
\toprule
信号组 & 包含 & 作用 \\
\midrule
命令/地址（CA） & CK, $\overline{CK}$, CKE, CS, RAS, CAS, WE, A[17:0], BA[1:0], BG[1:0] & 控制器→DRAM：发命令和地址 \\
数据（DQ） & DQ[63:0]（或 per-channel 8/16 bit） & 双向：读时 DRAM→控制器，写时控制器→DRAM \\
选通（DQS） & DQS[7:0], $\overline{DQS}$[7:0] & 数据选通：与 DQ 同源发出，接收端用 DQS 边沿采样 DQ \\
\bottomrule
\end{longtable}

DQS 是 DDR 接口的核心：它由发送端（读时是 DRAM，写时是控制器）与数据同源发出，频率与数据率相同。接收端用 DQS 的上升沿和下降沿分别采样 DQ，实现双倍数据率。

\begin{pdfdiagram}[x=1cm,y=1cm]
  % CK
  \draw[BookRule, line width=0.7pt] (0.5,3.0) -- (1.0,3.0) -- (1.0,3.6) -- (2.0,3.6) -- (2.0,3.0) -- (3.0,3.0) -- (3.0,3.6) -- (4.0,3.6) -- (4.0,3.0) -- (5.0,3.0) -- (5.0,3.6) -- (6.0,3.6) -- (6.0,3.0) -- (7.0,3.0);
  \node[hw label, anchor=east] at (0.4,3.3) {CK};
  % DQS
  \draw[BookAccent, line width=0.7pt] (0.5,1.8) -- (1.25,1.8) -- (1.25,2.4) -- (1.75,2.4) -- (1.75,1.8) -- (2.25,1.8) -- (2.25,2.4) -- (2.75,2.4) -- (2.75,1.8) -- (3.25,1.8) -- (3.25,2.4) -- (3.75,2.4) -- (3.75,1.8) -- (4.25,1.8) -- (4.25,2.4) -- (4.75,2.4) -- (4.75,1.8) -- (5.25,1.8) -- (5.25,2.4) -- (5.75,2.4) -- (5.75,1.8) -- (6.25,1.8) -- (6.25,2.4) -- (6.75,2.4) -- (6.75,1.8) -- (7.0,1.8);
  \node[hw label, anchor=east] at (0.4,2.1) {DQS};
  % DQ（眼图概念）
  \draw[BookTeal, line width=0.7pt] (0.5,0.6) -- (1.0,0.6) -- (1.25,1.0) -- (1.5,0.6) -- (1.75,1.0) -- (2.0,0.6) -- (2.25,1.0) -- (2.5,0.6) -- (2.75,1.0) -- (3.0,0.6) -- (3.25,1.0) -- (3.5,0.6) -- (3.75,1.0) -- (4.0,0.6) -- (4.25,1.0) -- (4.5,0.6) -- (4.75,1.0) -- (5.0,0.6) -- (5.25,1.0) -- (5.5,0.6) -- (5.75,1.0) -- (6.0,0.6) -- (6.25,1.0) -- (6.5,0.6) -- (6.75,1.0) -- (7.0,0.6);
  \node[hw label, anchor=east] at (0.4,0.8) {DQ};
  % 采样点
  \foreach \x in {1.25,1.75,2.25,2.75,3.25,3.75,4.25,4.75,5.25,5.75,6.25,6.75} {
    \fill[BookAmber] (\x,0.8) circle (0.06);
  }
  \node[BookAmber, font=\small\bfseries] at (7.5,0.8) {采样点};
  \node[hw label, text width=8cm] at (3.7,-0.5) {DQS 边沿（上升+下降）对齐 DQ 数据眼中心；接收端在 DQS 边沿采样 DQ。每个 DQS 周期采两次数据。};
\end{pdfdiagram}
\diagramnote{DDR 读时序：DQS 由 DRAM 与 DQ 同源发出，接收端（控制器 PHY）在 DQS 的每个边沿采样 DQ。DQS 与 DQ 之间的偏斜必须通过训练消除。}

\topic{PHY 训练：为什么必须做}

PCB 走线长度差异、过孔、连接器、封装引脚和 DRAM 内部延迟，使得 DQS 到达控制器 PHY 时与 DQ 的对齐关系不确定。训练的目标是调整 PHY 内部的延迟线，让采样点落在数据眼的正中心。

\begin{longtable}{L{0.24\linewidth}L{0.34\linewidth}L{0.32\linewidth}}
\toprule
训练步骤 & 做什么 & 调整对象 \\
\midrule
Write leveling & 控制器发 DQS，DRAM 回报 CK 与 DQS 的相位关系；调整 DQS 延迟使其对齐 CK & DQS 输出延迟 \\
Read gate training & 确定 DQS  preamble 的起始位置，防止 preamble 被误认为数据 & DQS 使能窗口 \\
Read leveling（per-bit） & 扫描 DQS 延迟，找到每根 DQ 的数据眼左右边界，取中心 & 每 bit 的 DQS 采样延迟 \\
Write leveling（per-bit） & 扫描 DQ 输出延迟，找到 DRAM 端数据眼中心 & 每 bit 的 DQ 输出延迟 \\
VREF training & 扫描接收端参考电压，找到眼图垂直中心 & PHY 的 VREF 值 \\
\bottomrule
\end{longtable}

训练在开机（cold boot）和从低功耗唤醒（warm boot）时执行。DDR5 还支持运行时周期性 retraining，补偿温度和电压漂移。

\topic{眼图}

把连续多个 bit 周期的 DQ 波形叠加在同一时间窗口，形成"眼图"。眼的张开程度反映信号质量：

\begin{longtable}{L{0.24\linewidth}L{0.34\linewidth}L{0.32\linewidth}}
\toprule
眼图参数 & 含义 & 典型要求（DDR4-3200） \\
\midrule
眼宽（水平） & 数据有效时间窗口 & $>$ 0.2 UI（约 125 ps） \\
眼高（垂直） & 噪声裕量 & $>$ 100 mV \\
抖动（jitter） & 边沿时间不确定度 & $<$ 5 ps RMS \\
串扰 & 相邻 DQ 线的耦合噪声 & 包含在眼高余量内 \\
\bottomrule
\end{longtable}

\topic{ODT 与驱动强度}

DDR 总线是多点拓扑（控制器 + 多颗 DRAM 挂在同一组信号上）。信号到达总线末端会反射，叠加在原始信号上破坏眼图。片上终端电阻（On-Die Termination, ODT）在 DRAM 内部接一个可编程电阻到 VDD 或 VSS，吸收到达的信号能量，抑制反射。

\begin{longtable}{L{0.22\linewidth}L{0.28\linewidth}L{0.4\linewidth}}
\toprule
ODT 设置 & 典型值 & 何时使用 \\
\midrule
40 $\Omega$ & 强终端 & 单 DIMM、短走线 \\
60 $\Omega$ & 中等终端 & 双 DIMM \\
120 $\Omega$ & 弱终端 & 远端 DRAM（减少负载） \\
Hi-Z & 关闭 & 该 DRAM 不是当前事务目标 \\
\bottomrule
\end{longtable}

驱动强度同理：控制器和 DRAM 的输出驱动器阻抗可编程（34/40/48 $\Omega$），与 ODT 配合形成匹配的传输线终端。

\section{四、内存控制器与调度}

\topic{控制器在系统中的位置}

内存控制器（Memory Controller, MC）位于 CPU die 内部（现代 x86 和 ARM SoC 均集成在片上），负责把 CPU 的 load/store 请求翻译成 DDR 命令序列，并管理刷新、ECC 和功耗状态。

\begin{pdfdiagram}[x=1cm,y=1cm]
  \node[hw state, minimum width=1.6cm] (cpu) at (0,0) {CPU 核};
  \node[hw compute, minimum width=1.8cm] (lsu) at (2.4,0) {LSU};
  \node[hw state, minimum width=1.6cm] (llc) at (4.8,0) {LLC};
  \node[hw control, minimum width=2.2cm] (mc) at (7.4,0) {内存控制器};
  \node[hw memory, minimum width=1.4cm] (phy) at (10.0,0) {PHY};
  \node[hw memory, minimum width=1.6cm] (dram) at (12.2,0) {DRAM};
  \draw[hw bus] (cpu) -- (lsu);
  \draw[hw bus] (lsu) -- (llc);
  \draw[hw bus] (llc) -- node[above, hw label]{miss} (mc);
  \draw[hw bus] (mc) -- (phy);
  \draw[hw bus] (phy) -- (dram);
  \node[hw label, text width=9cm] at (6.0,-1.3) {LLC miss 后请求进入内存控制器的请求队列；控制器按调度策略选择下一个发出的 DDR 命令。};
\end{pdfdiagram}
\diagramnote{从 CPU 到 DRAM 的路径。内存控制器是 LLC miss 和 DRAM 命令之间的调度层。}

\topic{请求队列与 bank 状态机}

控制器为每个 bank 维护一个状态机：

\begin{longtable}{L{0.18\linewidth}L{0.32\linewidth}L{0.4\linewidth}}
\toprule
Bank 状态 & 含义 & 可接受的命令 \\
\midrule
Idle（预充） & 行缓冲为空，无行打开 & ACTIVATE \\
Active（行打开） & 某行已在 row buffer 中 & READ / WRITE / PRECHARGE \\
Refreshing & 正在刷新，不可访问 & 无（等待完成） \\
Precharging & 正在关闭当前行 & 等待 tRP 后回到 Idle \\
\bottomrule
\end{longtable}

\topic{FR-FCFS 调度}

最常用的调度策略是 FR-FCFS（First-Ready, First-Come-First-Served）：

\begin{lstlisting}
priority 1: row-hit 请求（目标行已在 row buffer 中）→ 只需 CL
priority 2: row-empty 请求（bank idle）→ 需 tRCD + CL
priority 3: row-conflict 请求（需先 PRECHARGE）→ 需 tRP + tRCD + CL
tie-break:  同优先级内按到达顺序（FCFS）
\end{lstlisting}

Row-hit 请求代价最低（~14 ns），row-conflict 代价最高（~42 ns）。调度器优先发 row-hit 请求可以显著提高有效带宽。

\begin{longtable}{L{0.28\linewidth}L{0.22\linewidth}L{0.4\linewidth}}
\toprule
访问模式 & Row-hit 率 & 有效带宽（DDR4-3200, 25.6 GB/s 峰值） \\
\midrule
顺序流式 & $>$95\% & 22--24 GB/s \\
随机（4K 粒度） & 30--50\% & 8--14 GB/s \\
完全随机（64B） & $<$10\% & 3--6 GB/s \\
\bottomrule
\end{longtable}

\topic{读写切换代价}

DDR 总线是双向的：读时 DRAM 驱动 DQ，写时控制器驱动 DQ。切换方向需要一段死区（read-to-write turnaround 和 write-to-read turnaround），通常 2--4 个时钟周期。控制器把多个读请求和多个写请求分别攒成批次（batch），减少切换次数。

\topic{多通道交织}

现代 CPU 通常有 2--8 个内存通道。控制器把物理地址的某些位用于选择通道（interleave），使连续地址分布在不同通道上，多通道并行服务：

\begin{longtable}{L{0.22\linewidth}L{0.28\linewidth}L{0.4\linewidth}}
\toprule
配置 & 峰值带宽 & 交织粒度 \\
\midrule
单通道 DDR4-3200 & 25.6 GB/s & — \\
双通道 DDR4-3200 & 51.2 GB/s & 通常 64B 或 128B 粒度 \\
四通道 DDR5-4800 & 153.6 GB/s & 通常 64B 粒度 \\
八通道 DDR5-4800（服务器） & 307.2 GB/s & 通常 64B 粒度 \\
\bottomrule
\end{longtable}

\topic{ECC 与 patrol scrub}

服务器和高端桌面使用 ECC DRAM（每 64 数据位加 8 校验位，共 72 bit）。控制器在每次读回时校验并纠正单 bit 错误（SEC-DED）。Patrol scrub 是控制器周期性后台扫描全部内存，读出每行、校验、纠正并回写，防止单 bit 错误积累成不可纠正的多 bit 错误。

\section{五、NAND Flash 单元与 3D 堆叠}

\topic{浮栅单元：把电子困在绝缘层中}

传统 NAND Flash 单元是一个改造过的 MOSFET：在栅极和沟道之间插入一层被绝缘氧化物完全包围的导体（浮栅，floating gate）。电子通过隧穿进入浮栅后被绝缘层困住，即使断电也不会逃逸。浮栅中有电子 → 阈值电压升高 → 读时沟道不导通 → 代表 0；浮栅中无电子 → 阈值电压低 → 读时导通 → 代表 1。

\begin{pdfdiagram}[x=1cm,y=1cm]
  % 衬底
  \draw[fill=BookCodeBg] (0,-0.8) rectangle (6,0);
  \node[hw label] at (3,-0.4) {P 型衬底};
  % 隧穿氧化层
  \draw[fill=BookAmber!20] (1.5,0) rectangle (4.5,0.3);
  \node[hw label, font=\tiny] at (3,0.15) {隧穿氧化层 ($\sim$8 nm)};
  % 浮栅
  \draw[fill=BookAccent!30] (1.8,0.3) rectangle (4.2,0.7);
  \node[hw label, font=\tiny] at (3,0.5) {浮栅（多晶硅）};
  % 层间介质
  \draw[fill=BookAmber!20] (1.5,0.7) rectangle (4.5,1.0);
  \node[hw label, font=\tiny] at (3,0.85) {层间介质 (ONO)};
  % 控制栅
  \draw[fill=BookTeal!30] (1.8,1.0) rectangle (4.2,1.4);
  \node[hw label, font=\tiny] at (3,1.2) {控制栅（字线）};
  % 源漏
  \draw[fill=BookNavy!20] (0.2,-0.8) rectangle (1.2,0);
  \node[hw label, font=\tiny] at (0.7,-0.4) {N+源};
  \draw[fill=BookNavy!20] (4.8,-0.8) rectangle (5.8,0);
  \node[hw label, font=\tiny] at (5.3,-0.4) {N+漏};
  % 标注
  \node[hw label, text width=8cm] at (3,-1.6) {浮栅 NAND 单元截面。电子通过 FN 隧穿穿过 8nm 氧化层进入浮栅；被 ONO 层间介质困住，断电后保持 10 年以上。};
\end{pdfdiagram}
\diagramnote{浮栅 NAND 单元截面。控制栅（字线）施加高压时，电子隧穿进入浮栅（编程）；反向高压时电子隧穿离开（擦除）。}

\topic{电荷俘获（Charge Trap）单元}

3D NAND 普遍用氮化硅（SiN）电荷俘获层替代多晶硅浮栅。区别：浮栅是导体，电子在其中自由移动，一个缺陷点可能导致整层电荷泄漏；电荷俘获层是绝缘体，电子被局域陷阱捕获，单点缺陷只影响局部电荷，可靠性更好。

\begin{longtable}{L{0.2\linewidth}L{0.35\linewidth}L{0.35\linewidth}}
\toprule
特性 & 浮栅（Floating Gate） & 电荷俘获（Charge Trap） \\
\midrule
存储层材料 & 多晶硅（导体） & SiN（绝缘体） \\
电子状态 & 自由移动 & 局域陷阱 \\
缺陷影响 & 单点缺陷→整层泄漏 & 单点缺陷→局部影响 \\
单元间耦合 & 有（相邻浮栅电容耦合） & 极小 \\
适合 3D 堆叠 & 困难（耦合干扰） & 适合（垂直沟道） \\
\bottomrule
\end{longtable}

\topic{SLC/MLC/TLC/QLC：每单元存多少 bit}

通过精确控制浮栅/陷阱中的电子数量，可以把阈值电压分成多个区间，每个区间代表一个状态：

\begin{longtable}{L{0.12\linewidth}L{0.14\linewidth}L{0.16\linewidth}L{0.16\linewidth}L{0.32\linewidth}}
\toprule
类型 & 状态数 & bit/单元 & 阈值电压窗口 & P/E 寿命 \\
\midrule
SLC & 2 & 1 & $\sim$5V 全窗口 & $\sim$100,000 次 \\
MLC & 4 & 2 & 每级 $\sim$2V & $\sim$3,000--10,000 次 \\
TLC & 8 & 3 & 每级 $\sim$0.8V & $\sim$1,000--3,000 次 \\
QLC & 16 & 4 & 每级 $\sim$0.4V & $\sim$500--1,000 次 \\
\bottomrule
\end{longtable}

级数越多，相邻状态的阈值电压间隔越窄，对噪声和电荷泄漏越敏感，读写越慢，寿命越短。

\topic{编程算法：增量步进脉冲（ISPP）}

NAND 编程不是一次完成的。控制器施加一系列递增的编程脉冲（每次电压增加 $\Delta V_{ISPP}$，典型 0.2--0.5V），每施加一个脉冲后验证（read-verify）单元阈值电压是否已达到目标区间。未达到则继续下一个脉冲。

\begin{longtable}{L{0.28\linewidth}L{0.22\linewidth}L{0.4\linewidth}}
\toprule
参数 & TLC 典型值 & 含义 \\
\midrule
ISPP 步进 & 0.2--0.4 V & 每个编程脉冲增加的栅压 \\
验证次数 & 3--7 次/页 & 达到目标状态所需的脉冲数 \\
页编程时间 & 0.5--2 ms & 一整页（16 KiB）编程完成 \\
读干扰 & 每读 $\sim$100K 次需刷新 & 未选中单元的栅极电压累积效应 \\
\bottomrule
\end{longtable}

\topic{3D NAND：垂直堆叠}

2D NAND 的面积缩小已接近极限（单元间耦合干扰）。3D NAND 把字线层垂直堆叠（目前量产 232--300+ 层），沟道变成垂直的圆柱体穿过所有字线层：

\begin{pdfdiagram}[x=1cm,y=1cm]
  % 垂直沟道
  \draw[fill=BookTeal!20] (2.8,-2.0) rectangle (3.2,2.5);
  \node[hw label, rotate=90] at (3.0,0.3) {垂直沟道};
  % 字线层
  \foreach \y/\n in {-1.8/WL0, -1.2/WL1, -0.6/WL2, 0/WL3, 0.6/WL4, 1.2/WL5, 1.8/WL6} {
    \draw[fill=BookAccent!25] (0.5,\y) rectangle (2.6,\y+0.35);
    \draw[fill=BookAccent!25] (3.4,\y) rectangle (5.5,\y+0.35);
    \node[hw label, font=\tiny] at (-0.1,\y+0.17) {\n};
  }
  % 层间介质
  \node[hw label, font=\tiny] at (4.5,2.3) {层间介质};
  % 标注
  \node[hw label, text width=8cm] at (3.0,-2.8) {3D NAND 截面：垂直沟道（多晶硅柱）穿过数十到数百层字线。每层字线与沟道交叉处形成一个存储单元。电荷俘获层包裹在沟道周围。};
\end{pdfdiagram}
\diagramnote{3D NAND 结构。232 层意味着一根垂直沟道上串联 232 个存储单元。多个沟道柱组成一个 string，多个 string 组成一个 block。}

\begin{longtable}{L{0.28\linewidth}L{0.22\linewidth}L{0.4\linewidth}}
\toprule
参数 & 典型值（232 层） & 含义 \\
\midrule
字线层数 & 232 & 每根沟道上的单元数 \\
String 长度 & 232 单元串联 & 读时电流路径长，速度慢 \\
Die 容量 & 1--2 Tbit & 单 die \\
Page 大小 & 16--36 KiB & 一次编程的最小单位 \\
Block 大小 & 4--16 MiB & 一次擦除的最小单位 \\
\bottomrule
\end{longtable}

\topic{Multi-plane 与 die 级并行}

一颗 NAND die 内部有多个 plane（通常 2--4 个），每个 plane 有独立的 page register 和行译码。Multi-plane 操作允许同时对多个 plane 发相同的命令（如同时读 4 个 plane 的同一页），把内部带宽提高 2--4 倍。

\begin{longtable}{L{0.24\linewidth}L{0.26\linewidth}L{0.4\linewidth}}
\toprule
并行层级 & 粒度 & 带宽提升 \\
\midrule
Plane 内 & 页（16 KiB） & 基础单位 \\
Multi-plane & 2--4 plane 同时操作 & 2--4$\times$ \\
Multi-die（多 LUN） & 2--4 die 同时操作 & 再 2--4$\times$ \\
Multi-channel（控制器） & 4--8 通道 & 再 4--8$\times$ \\
\bottomrule
\end{longtable}

\topic{FTL 映射与 GC 回顾}

第五章已介绍 FTL 的基本映射和写放大。这里补充结构细节：

\begin{longtable}{L{0.24\linewidth}L{0.34\linewidth}L{0.32\linewidth}}
\toprule
FTL 组件 & 作用 & 存储位置 \\
\midrule
L2P 映射表 & 逻辑页号→物理页号 & DRAM 或 HMB（主机内存缓冲） \\
有效页位图 & 每 block 中哪些页仍有效 & SRAM / DRAM \\
坏块表 & 出厂和使用中产生的坏块 & 独立区域 \\
GC 候选队列 & 有效页最少的 block 优先回收 & 控制器 SRAM \\
掉电保护日志 & 映射表更新的事务日志 & NAND 专用区域 + 电容 \\
\bottomrule
\end{longtable}

\section{六、磁盘界面物理与伺服}

\topic{空气轴承：磁头怎样飞起来}

硬盘磁头不接触盘片。盘片高速旋转时，表面带动一层空气形成气流；磁头滑块（slider）底面刻有精密的导轨几何（rail），气流在导轨下方产生正压把滑块抬起，同时在导轨后缘产生负压防止滑块飞得太高。两个力平衡后，磁头稳定悬浮在盘片上方 5--10 nm 处。

\begin{pdfdiagram}[x=1cm,y=1cm]
  % 盘片表面
  \draw[BookRule, line width=0.7pt] (0,0) -- (10,0);
  \node[hw label] at (5,-0.4) {盘片表面（旋转方向 $\to$）};
  % 气流
  \draw[BookTeal, line width=0.5pt, ->] (0.5,0.1) -- (2.0,0.1);
  \draw[BookTeal, line width=0.5pt, ->] (0.5,0.3) -- (2.0,0.3);
  \draw[BookTeal, line width=0.5pt, ->] (0.5,0.5) -- (2.0,0.5);
  % 滑块
  \draw[fill=BookCodeBg, line width=0.6pt] (3.0,0.15) -- (3.0,0.8) -- (7.0,0.8) -- (7.0,0.15) -- (6.5,0.08) -- (5.5,0.12) -- (4.5,0.08) -- (3.5,0.12) -- cycle;
  \node[hw label, font=\tiny] at (5.0,0.5) {滑块（slider）};
  % 磁头
  \draw[fill=BookAccent!40] (6.2,0.08) rectangle (6.8,0.3);
  \node[hw label, font=\tiny] at (6.5,-0.15) {磁头};
  % 飞行高度
  \draw[BookAmber, line width=0.5pt, <->] (7.5,0) -- (7.5,0.12);
  \node[BookAmber, font=\tiny] at (8.2,0.06) {5--10 nm};
  % 正压/负压标注
  \node[hw label, font=\tiny] at (4.0,1.1) {正压区（抬起）};
  \node[hw label, font=\tiny] at (6.0,1.1) {负压区（下压）};
  \node[hw label, text width=8cm] at (5.0,-1.2) {空气轴承原理：盘片旋转带动气流，滑块底面导轨几何产生正压（前缘）和负压（后缘），平衡后磁头悬浮在 5--10 nm 高度。};
\end{pdfdiagram}
\diagramnote{磁头飞行原理。飞行高度由滑块几何、转速和空气粘度共同决定。转速越高飞行越高；盘片倾斜或振动会改变局部飞行高度。}

\topic{热飞行高度控制（TFC）}

现代硬盘在磁头附近集成一个微型加热器。通电后加热器使滑块局部热膨胀，磁头向盘片靠近 1--2 nm。控制器通过调节加热功率精确控制飞行高度：读写时需要更近（信号更强），空闲时拉远（减少磨损风险）。

\topic{伺服系统：磁头怎样找到目标磁道}

盘片上除了数据扇区，还嵌有伺服扇区（servo sector），每个伺服扇区包含：

\begin{longtable}{L{0.22\linewidth}L{0.34\linewidth}L{0.34\linewidth}}
\toprule
伺服字段 & 内容 & 作用 \\
\midrule
Preamble & 固定频率图案 & 自动增益控制（AGC）和时钟恢复 \\
Servo Address Mark (SAM) & 特殊图案 & 标识伺服扇区起始 \\
Gray code & 磁道号（格雷码） & 粗定位：知道当前在哪条磁道 \\
Burst（A/B/C/D） & 偏移图案 & 精定位：计算磁头相对磁道中心的偏移量（PES） \\
\bottomrule
\end{longtable}

磁头每经过一个伺服扇区（典型每秒约 20,000--40,000 次），读出 Gray code 和 burst，计算位置误差信号（Position Error Signal, PES）。音圈电机（Voice Coil Motor, VCM）根据 PES 驱动磁头臂，形成闭环：

\begin{lstlisting}
servo loop (runs at servo rate ~25 kHz):
  read burst A/B/C/D from servo sector
  compute PES = f(A, B, C, D)   // 位置误差
  PID controller: VCM_current = Kp*PES + Ki*integral + Kd*derivative
  actuator moves head toward track center
\end{lstlisting}

\begin{longtable}{L{0.28\linewidth}L{0.22\linewidth}L{0.4\linewidth}}
\toprule
参数 & 典型值 & 含义 \\
\midrule
伺服采样率 & 25--40 kHz & 每秒读伺服扇区的次数 \\
磁道密度 & 200K--300K TPI & 每英寸磁道数；磁道间距 $\sim$80--130 nm \\
定位精度（3$\sigma$） & $<$10\% 磁道间距 & 约 8--13 nm \\
寻道时间（平均） & 4--9 ms & 音圈加速+减速+稳定 \\
\bottomrule
\end{longtable}

\topic{磁记录物理：超顺磁极限与 HAMR}

磁记录把 bit 写成一小块磁性薄膜的磁化方向。每个 bit 由数百个磁颗粒（grain）组成。颗粒越小，热扰动越容易翻转磁化方向（超顺磁效应），数据越不稳定。这是面密度增长的根本物理限制。

\begin{longtable}{L{0.18\linewidth}L{0.28\linewidth}L{0.44\linewidth}}
\toprule
技术 & 原理 & 面密度 \\
\midrule
LMR（纵磁） & 磁化方向平行于盘面 & $\sim$100 Gb/in$^2$（已淘汰） \\
PMR（垂磁） & 磁化方向垂直于盘面 & $\sim$1 Tb/in$^2$ \\
SMR（叠瓦） & 磁道部分重叠，写时重写相邻带 & $\sim$1.2 Tb/in$^2$ \\
HAMR（热辅助） & 激光局部加热介质降低矫顽力后写入 & $>$2 Tb/in$^2$（量产中） \\
MAMR（微波辅助） & 微波场辅助翻转 & 研发中 \\
\bottomrule
\end{longtable}

HAMR 在磁头中集成一个激光二极管，写入前用激光把介质局部加热到 400--500°C（矫顽力急剧下降），写入后介质迅速冷却，磁化方向被"冻结"。这允许使用矫顽力极高（热稳定性好）的介质，突破超顺磁极限。

\topic{盘片振动与 off-track 容错}

旋转盘片会受到外部振动、气流扰动和主轴电机抖动的影响，导致盘片相对磁头产生微小位移。控制器通过以下机制应对：

\begin{longtable}{L{0.28\linewidth}L{0.34\linewidth}L{0.28\linewidth}}
\toprule
机制 & 作用 & 典型参数 \\
\midrule
伺服闭环带宽 & 跟踪低频振动 & 1--3 kHz \\
前馈加速度计 & 检测外部冲击，提前补偿 & 双轴 MEMS 加速度计 \\
重复性跳动（RRO）补偿 & 学习每转一次的周期性偏移并前馈消除 & 存储每磁道的偏移表 \\
非重复性跳动（NRRO） & 随机振动，只能靠闭环带宽抑制 & 3$\sigma$ $<$ 5 nm \\
\bottomrule
\end{longtable}

当 off-track 量超过磁道间距的 $\sim$15\% 时，读信号质量急剧下降，控制器触发重试或重新寻道。

\section{七、存储层次总览}

把本章所有器件放在同一张表中，可以看到物理结构如何决定系统参数：

\begin{longtable}{L{0.12\linewidth}L{0.14\linewidth}L{0.12\linewidth}L{0.12\linewidth}L{0.12\linewidth}L{0.14\linewidth}L{0.14\linewidth}}
\toprule
器件 & 物理机制 & 容量 & 读延迟 & 写延迟 & 保持 & 失效模式 \\
\midrule
寄存器 & 触发器 & KB & $<$1 ns & $<$1 ns & 通电 & 断电丢失 \\
SRAM (L1/L2) & 6T 锁存 & 32K--16M & 1--4 ns & 1--4 ns & 通电 & 断电丢失；SNM 不足→软错误 \\
DRAM & 1T1C 电容 & 8--128 GB & 50--100 ns & 50--100 ns & 需刷新 & 漏电→数据丢失；刷新失败 \\
NAND SSD & 浮栅/电荷俘获 & 256G--8T & 25--100 $\mu$s & 200 $\mu$s--2 ms & 10 年 & P/E 耗尽→坏块；读干扰 \\
HDD & 磁畴方向 & 1--24 TB & 4--12 ms & 4--12 ms & 数十年 & 磁头碰撞；伺服丢失；介质退化 \\
\bottomrule
\end{longtable}

每一层的访问粒度、写入代价和失效模式都由其物理结构决定，不能通过软件抽象消除。理解这些结构，才能理解为什么 cache miss 代价是 100 ns 而不是 1 ns，为什么 SSD 随机写比顺序写慢，为什么 HDD 随机 IOPS 只有 100 而顺序带宽有 200 MB/s。
